% SLIDE 7: CSTM --- Image Space
\begin{frame}{Canonical Camera Space Transform --- Image Space}
\begin{columns}[T]
\column{0.55\textwidth}
\textbf{CSTM\textsubscript{image}} resizes the input image instead of scaling labels:
\[
  \boxed{I_c = \mathcal{T}(I,\; \omega_r), \qquad \omega_r = \frac{f^c}{f}}
\]

\textbf{Why resize?} After resizing, objects in $I_c$ appear at the angular size they would subtend through the canonical camera. The encoder sees a consistent ``canonical viewpoint.''

\vspace{0.3cm}
\textbf{Both paths are equivalent:}
\begin{itemize}
  \item CSTM\textsubscript{label}: scale depth labels, keep images
  \item CSTM\textsubscript{image}: resize images, keep depth labels
\end{itemize}

In practice, the paper uses \textbf{CSTM\textsubscript{image}} during training.

\column{0.42\textwidth}
\centering
\begin{tikzpicture}[scale=0.75]
  % Original image (wide-angle)
  \draw[thick, atlantic, fill=atlantic!5] (0,3) rectangle (3.5,5.5);
  \node[font=\scriptsize\sffamily, color=atlantic] at (1.75,5.8) {Wide-angle: $f = 500$};
  \node[font=\scriptsize, color=bk70] at (1.75,4.25) {Small objects};

  % Arrow
  \draw[dataarrow=garnet, very thick] (4.0,4.25) -- (5.0,4.25)
    node[midway, above, font=\scriptsize\sffamily] {$\omega_r = 2$};

  % Canonical image (zoomed)
  \draw[thick, garnet, fill=garnet!5] (5.5,2.5) rectangle (10.0,6.0);
  \node[font=\scriptsize\sffamily, color=garnet] at (7.75,6.3) {Canonical: $f^c = 1000$};
  \node[font=\scriptsize, color=bk70] at (7.75,4.25) {Enlarged objects};

  % Crop indicator
  \draw[dashed, rose, thick] (6.0,3.0) rectangle (9.5,5.5);
  \node[font=\scriptsize, color=rose] at (7.75,2.6) {Crop to input size};
\end{tikzpicture}
\end{columns}
\end{frame}
